\documentclass[a4paper,12pt]{extarticle}
\usepackage[table,xcdraw]{xcolor}
\usepackage{pgfplots}
\usepackage[unicode, pdftex]{hyperref}
\usepackage[T2A]{fontenc}
\usepackage[utf8]{inputenc}
\usepackage[english,russian]{babel}
\usepackage{amsmath,amsfonts,amsthm, mathtools}
\usepackage{amssymb}
\usepackage{dsfont}
\usepackage{graphicx}
\graphicspath{{images/}}
\usepackage{wrapfig}
\usepackage{floatrow}
\usepackage{setspace}
\usepackage[headings]{fancyhdr}
\usepackage[margin=10pt,font={small,stretch=0,9},labelfont=bf,labelsep=period,
justification=centerlast]{caption}
\usepackage{array,tabularx,tabulary,booktabs}
\RequirePackage{multirow}
\RequirePackage{multicol}
\RequirePackage{longtable}
\usepackage{indentfirst}
\usepackage{hyperref}
\usepackage{icomma}
\newcommand*{\hm}[1]{#1\nobreak\discretionary{}
	{\hbox{$\mathsurround=0pt #1$}}{}}
\usepackage{soulutf8} 
\usepackage{etoolbox} 
\usepackage{geometry}
\geometry{top=20mm}
\geometry{bottom=20mm}
\geometry{left=20mm}
\geometry{right=20mm}
\DeclareMathOperator{\sgn}{\mathop{sgn}}
\usepackage[OT1]{fontenc}
\usepackage{amsmath}
\usepackage{amsfonts}
\usepackage{amssymb}
\usepackage{wasysym}
\usepackage{wrapfig}
\usepackage{physics}
\usepackage[normalem]{ulem}
\graphicspath{{pictures/}}
\DeclareGraphicsExtensions{.png,.jpg,.jpeg}
\usepackage{indentfirst}
\usepackage[T2A]{fontenc}
\usepackage{subfigure}
\usepackage{fancyhdr}
\usepackage{caption}

\usepackage{tabularx}
\usepackage{floatrow}
\usepackage{multicol}
\usepackage{lastpage}

\pgfplotsset{
  MyStyle1/.style={
    axis lines = middle,
    enlargelimits = true,
    grid=both,
    grid style={line width=.1pt, draw=gray!10},
    major grid style={line width=.2pt,draw=gray!50},
    minor tick num=5,
    enlargelimits=0.05,
    ticklabel style={font=\tiny,fill=white},
    xlabel style={at={(ticklabel* cs:1)},anchor=north west},
    ylabel style={at={(ticklabel* cs:1)},anchor=south east},
  }
}
\def\Slope{1.5}
\def\Intercept{-20}


% Настройка отчёркиваний и нумерации
\usepackage{fancyhdr}
\pagestyle{fancy}
\fancyhf{}

\fancyhead[L]{\nouppercase{\leftmark\hfill}}
\fancyfoot[RO, RE]{Страница \thepage \ из \pageref{LastPage}}
\renewcommand{\headrulewidth}{0.5pt}
\renewcommand{\footrulewidth}{0.5pt}

\usepackage{titlesec}
\titlelabel{\thetitle.\quad}

\pgfplotsset{compat=1.18}
\begin{document}

\thispagestyle{empty}
\begin{center}
	\textit{Федеральное государственное автономное образовательное\\ учреждение высшего образования }
	
	\vspace{0.5ex}
	
	\textbf{«Московский физико-технический институт\\ (национальный исследовательский университет)»}
\end{center}

\vspace{10ex}
\begin{figure}[!h]
  \begin{center}
    
\includegraphics[width = 0.8\textwidth]{MIPT.png}
  \end{center}
\end{figure}

\begin{center}
	
	\textbf{Лабораторная работа № 4.4.4}
	
	\vspace{1ex}
	
	по курсу общей физики
	
	на тему:
	
	\textbf{\textit{<<Интерферометр Фабри-Перо>>}}
	
	\vspace{20ex}
	
	\begin{flushright}
		\noindent
		\textit{Работу выполнил:}\\  
		\textit{Легких Никита \\(группа Б05-301)}
	\end{flushright}
	\vfill
	Долгопрудный \\ \today
	\setcounter{page}{1}
\end{center}
\newpage

% Конец шаблона
\newpage

\section*{Введение}
\textbf{Цель работы:} определение спектральных характеристик интерферометра Фабри-Перо на основании наблюдаемых дифракционных колец ртутной лампы.\\
\textbf{В работе используются:}  интерферометр Фабри-Перо, линзы, светофильтры, ртутная лампа, катетометр КН-6.
\section*{Экспериментальная установка}
\begin{figure}[H]
	\centering
	\includegraphics[width=0.86\textwidth]{pics/ustanovka.png}
    \caption{Схема экспериментальной установки}
\end{figure}
На схеме $S$ --- лампа, $ Л_0 $ --- линза, $ C $ --- светофильтр, ИФП --- интерферометр
Фабри-Перо, $T$ --- зрительная труба. Диаметры колец измеряются с помощью микроскопа
катетометра.\\
Цена деления катетометра КН-6 составляет $0,001$ мм, однако в лабораторной работе мы измеряли с точностью до десятых, так же, в силу погрешности при измерении середины кольца, примем $\Delta = 0,01$ мм.
\section*{Ход работы}
Проведём требуемые измерения на ртутной лампе, меняя светофильтр и записывая координаты соответствующих колец. Всего удалось измерить зелёную и две жёлтые линии, по пять измерений на каждую. Полученные данные внесём в таблицы ниже:
\begin{table}[H]
\centering
\begin{tabular}{|c|c|c|}
\hline
$x^{\text{up}}$, мм & $x^{\text{down}}$, мм & $d$, мм \\ \hline
199,86 & 165,04 & 34,82 \\ \hline
198,11 & 166,86 & 31,25 \\ \hline
195,95 & 169,12 & 26,83 \\ \hline
193,52 & 171,28 & 22,24 \\ \hline
190,22 & 175,03 & 15,19 \\ \hline
\end{tabular}
\caption{Измерения для зелёной линии}
\end{table}
\begin{table}[H]
\centering
\begin{tabular}{|c|c|c|c|c|c|}
\hline
$x^{\text{up}}$, мм & $x^{\text{down}}$, мм & $d$, мм & $x^{\text{up}}$, мм & $x^{\text{down}}$, мм & $d$, мм\\ \hline
200,56 & 164,40 & 36,16 & 199,92 & 165,20 & 34,72 \\ \hline
198,74 & 166,09 & 32,65 & 198,00 & 167,02 & 30,98 \\ \hline
196,55 & 168,45 & 28,10 & 195,75 & 169,31 & 26,44 \\ \hline
194,00 & 171,00 & 23,00 & 193,00 & 172,09 & 20,91 \\ \hline
191,04 & 174,26 & 16,78 & 189,20 & 175,78 & 13,42 \\ \hline
\end{tabular}
\caption{Измерения для жёлтых линий}
\end{table}
Погрешность прямых измерений ($x^{\text{up}}$ и $x^{\text{down}}$) совпадает с погрешностью катетометра $\sigma = 0,01$ мм, погрешность же $d = x^{\text{up}}-x^{\text{down}}$ равна $\sigma_d = 0,02$ мм.\\
Построим зависимость квадрата диаметра от номера кольца $d^2(n)$. Для этого проведем предварительную обработку данных:
\begin{table}[H]
\centering
\begin{tabular}{|c|cc|cc|cc|}
\hline
  & \multicolumn{2}{c|}{Зел. ($\lambda = 546,07$ нм)}                  & \multicolumn{2}{c|}{Жел. ($\lambda = 579,07$ нм)}               & \multicolumn{2}{c|}{Жел. ($\lambda = 576,96$ нм)}               \\ \hline
$n$ & \multicolumn{1}{c|}{$d^2$}       & $\sigma_{d^2}$ & \multicolumn{1}{c|}{$d^2$}       & $\sigma_{d^2}$ & \multicolumn{1}{c|}{$d^2$}       & $\sigma_{d^2}$ \\ \hline
1 & \multicolumn{1}{c|}{230,74}  & 0,43  & \multicolumn{1}{c|}{281,57}  & 0,47  & \multicolumn{1}{c|}{180,10}  & 0,38  \\ \hline
2 & \multicolumn{1}{c|}{494,62}  & 0,63  & \multicolumn{1}{c|}{529,00}  & 0,65  & \multicolumn{1}{c|}{437,23}  & 0,59  \\ \hline
3 & \multicolumn{1}{c|}{719,85}  & 0,76  & \multicolumn{1}{c|}{789,61}  & 0,79  & \multicolumn{1}{c|}{699,07}  & 0,75  \\ \hline
4 & \multicolumn{1}{c|}{976,56}  & 0,88  & \multicolumn{1}{c|}{1066,02} & 0,92  & \multicolumn{1}{c|}{959,76}  & 0,88  \\ \hline
5 & \multicolumn{1}{c|}{1212,43} & 0,98  & \multicolumn{1}{c|}{1307,55} & 1,02  & \multicolumn{1}{c|}{1205,48} & 0,98  \\ \hline
\end{tabular}
\caption{Обработка данных}
\end{table}
На основании этого построим соответствующие графики:

График для $\lambda = 546.07$ нм:
\begin{minipage}{0.6\textwidth}
      \begin{tikzpicture}
        \begin{axis}[
          width=1\textwidth,
          height=,
          xlabel={$n$},
          ylabel={$d^2,$ мм$^2$},
          MyStyle1,
          ]
          \draw[blue, domain=0:20] plot (\x,244.48000000000005*\x-6.5480000000001155);
          \addplot [only marks] plot [error bars/.cd,
                x dir=both, x fixed = 0.05,
                y dir=both, y explicit,
                      ] table [
                      x=n, y=1d,
                x error minus=null, x error plus=null, y error minus=delta1d, y error plus=delta1d, col sep=comma] {all.csv};
        \end{axis}
      \end{tikzpicture}
\end{minipage}

График для $\lambda = 549.07$ нм:
\begin{minipage}{0.6\textwidth}
      \begin{tikzpicture}
        \begin{axis}[
          width=1\textwidth,
          height=,
          xlabel={$n$},
          ylabel={$d^2,$ мм$^2$},
          MyStyle1,
          ]
          \draw[blue, domain=0:20] plot (\x,258.8980000000001*\x+18.055999999999813);
          \addplot [only marks] plot [error bars/.cd,
                x dir=both, x fixed = 0.05,
                y dir=both, y explicit,
                      ] table [
                      x=n, y=2d,
                x error minus=null, x error plus=null, y error minus=delta2d, y error plus=delta2d, col sep=comma] {all.csv};
        \end{axis}
      \end{tikzpicture}
\end{minipage}

График для $\lambda = 576.96$ нм:
\begin{minipage}{0.6\textwidth}
      \begin{tikzpicture}
        \begin{axis}[
          width=1\textwidth,
          height=,
          xlabel={$n$},
          ylabel={$d^2,$ мм$^2$},
          MyStyle1,
          ]
          \draw[blue, domain=0:20] plot (\x,257.329*\x-75.6590000000001);
          \addplot [only marks] plot [error bars/.cd,
                x dir=both, x fixed = 0.05,
                y dir=both, y explicit,
                      ] table [
                      x=n, y=1d,
                x error minus=null, x error plus=null, y error minus=delta1d, y error plus=delta1d, col sep=comma] {all.csv};
        \end{axis}
      \end{tikzpicture}
\end{minipage}

С помощью МНК найдем угловые коэффициенты соответствующих зависимостей и, исходя из них, по формулам: $\frac{\lambda}{L} = \frac{k}{4 f^2}$ и $\Delta L =\frac{\Delta k}{k} \cdot L$, где $f = 110$ мм - фокусное расстояние линзы, оценим базу интерферометра $L$:\\
$k_1 = (247,9 \pm 2,3)$ $\text{мм}^2$ $\Rightarrow$ $L_1 = (0,1081 \pm 0,0014)$ мм\\
$k_2 = (258,9 \pm 3)$ $\text{мм}^2$ $\Rightarrow$ $L_2 = (0,1083 \pm 0,0013)$ мм\\
$k_3 = (257,3 \pm 1,6)$ $\text{мм}^2$ $\Rightarrow$ $L_3 = (0,1085 \pm 0,0007)$ мм\\
Как видно, значения $L$ сходятся с очень хорошей точностью в пределах доверительных интервалов. Усреднив эти результаты, получим $L = (0,1083 \pm 0,0011)$ мм, что достаточно близко к теоретическому значению $L_{\text{теор}} = 0,1$ мм.\\
Также посчитаем $d_{\text{ср}}$ для желтого дублета и построим график зависимости $d_{\text{ср}}(1/\Delta d)$:
\begin{table}[H]
\centering
\begin{tabular}{|c|c|c|c|}
\hline
$n$ & $d_{\text{ср}}$, мм & $1/\Delta d$, $\text{мм}^{-1}$ & $\sigma_{1/\Delta d}$, $\text{мм}^{-1}$ \\ \hline
1 & 15,100 & 0,298 & 0,004 \\ \hline
2 & 21,955 & 0,478 & 0,009 \\ \hline
3 & 27,270 & 0,602 & 0,015 \\ \hline
4 & 31,815 & 0,599 & 0,014 \\ \hline
5 & 35,440 & 0,694 & 0,019 \\ \hline
\end{tabular}
\caption{Обработка данных для жёлтого дублета}
\end{table}
График для желтого дуплета:
\begin{minipage}{0.6\textwidth}
      \begin{tikzpicture}
        \begin{axis}[
          width=1\textwidth,
          height=,
          xlabel={$\frac{1}{\Delta d}$, мм$^{-1}$},
          ylabel={$d_{\text{ср}},$ мм},
          MyStyle1,
          ]
          \draw[blue, domain=0:3] plot (\x,51.17172022538402*\x-1.01993294440015);
          \addplot [only marks] plot [error bars/.cd,
                x dir=both, x fixed = 0.05,
                y dir=both, y explicit,
                      ] table [
                      x=1_d, y=d,
                x error minus=delta1_d, x error plus=delta1_d, y error minus=null, y error plus=null, col sep=comma] {duplet.csv};
        \end{axis}
      \end{tikzpicture}
\end{minipage}

С помощью МНК найдем угловой коэффициент зависимости и, исходя из него, по формуле $k = \frac{4 f^2 \Delta \lambda}{\lambda}$ оценим разность длины волн для жёлтого дублета $\Delta \lambda$:\\
$k = (52,2 \pm 7,3)$ $\Rightarrow$ $\Delta \lambda = (0,62 \pm 0,14)$ нм\\
Как видно, это значение близко к теоретическому: $\Delta \lambda_{\text{теор}} = 2,11$ нм по порядку величины, однако всё же не сходится с ним в пределах доверительного интервала. Также стоит отметить, что точность данного метода низка - $\varepsilon = 23 \%$. Таким образом, данный метод подходит скорее для оценки порядка величины, нежели для получения точного значения.\\
На основании полученных значений оценим экспериментальное и теоретическое значения линейной дисперсии интерферометра с помощью следующих формул:
\begin{equation*}
    D_{\text{эксп}} = \frac{\Delta D}{2 \Delta \lambda} \text{ мм/}{\text{\AA}}
\end{equation*}
\begin{equation*}
    D_{\text{теор}} = \frac{2f^2}{\lambda D} \text{ мм/}{\text{\AA}}
\end{equation*}
Полученные данные внесем в таблицу ниже:
\begin{table}[H]
\centering
\begin{tabular}{|c|c|c|c|}
\hline
$D_{\text{эксп}}$, \text{мм/}{\text{\AA}} & $\sigma_{D_{\text{эксп}}}$, \text{мм/}{\text{\AA}} & $D_{\text{теор}}$, \text{мм/}{\text{\AA}} & $\sigma_{D_{\text{теор}}}$, \text{мм/}{\text{\AA}} \\ \hline 
0,27 & 0,05 & 0,277 & 0,003 \\ \hline
0,17 & 0,03 & 0,190 & 0,004 \\ \hline
0,13 & 0,03 & 0,153 & 0,004 \\ \hline
0,13 & 0,03 & 0,131 & 0,003 \\ \hline
0,12 & 0,02 & 0,118 & 0,003 \\ \hline
\end{tabular}
\caption{Линейная дисперсия}
\end{table}
Как видно, для всех точек теоретическое и экспериментальное значения равны в пределах доверительных интервалов.\\
Также измерим толщину зелёного кольца:
\begin{align*}
\delta r = (0,71 \pm 0,02)\ \text{мм}
\end{align*}
Из этих данных получим аппаратную разрешающую способность:  
\begin{align*}
  &R = \dfrac{\lambda}{\delta \lambda} \approx \dfrac{4f^2}{d \delta r} = (3100 \pm 100) \quad \text{для зелёной линии}
\end{align*}
Оценим также прочие характеристики интерферометра:\\
- максимальный порядок интерференции $m$ для желтых линий ртути:\\
$$ m = \frac{2L \cos \theta}{\lambda} \approx \frac{2L}{\lambda} = (187,0 \pm 1,9)$$
- добротность (при $r = 0.0913$):
$$ Q = \frac{2 \pi L}{\lambda (1 - r)} = (4700 \pm 50)$$
- число интерферирующих лучей: 
$$ N = \frac{Q}{m} = (25,1 \pm 0,4)$$
\section*{Выводы}
В ходе работы удалось провести оценку некоторых характеристик интерферометра:
\begin{itemize}
    \item Была оценена база интерферометра $L = (0,1083 \pm 0,0011)$ мм, что достаточно близко к теоретическому значению $L_{\text{теор}} = 0,1$ мм.
    \item Для жёлтого дублета была найдена разность длин волн $\Delta \lambda = (0,62 \pm 0,14)$ нм, что сходится с теоретическим значением $\Delta \lambda_{\text{теор}} = 2,11$ по порядку величины, однако не в рамках доверительного интервала, к тому же, точность данного метода достаточно мала - $\varepsilon = 23 \%$.
    \item Были посчитаны значения линейной дисперсии интерферометра, значения сошлись с теоретическими в пределах доверительных интервалов.
    \item Были оценены такие параметры как: аппаратная разрешающая способность - $R = (3100 \pm 100)$, максимальный порядок интерференции - $m = (187,0 \pm 1,9)$, добротность - $Q = (4700 \pm 50)$, число интерферирующих лучей - $N = (25,1 \pm 0,4)$.
\end{itemize}

\end{document}